\section{Source-Grounded Reanalysis and Intervention Details}
\subsection{Count reconstruction}
The financial audit reports 101 success-derived retrievals at 0.931 accuracy and 34 failure-derived retrievals at 0.647. The unique integer numerators consistent with three-decimal rounding are 94 and 22, yielding 7 and 12 failures. It also reports ten W$\rightarrow$W persistent failures, all failure-derived; hence 10/12 failure-derived errors are persistent, while the success-derived stratum contains no W$\rightarrow$W case.

\subsection{Terminal information-equivalence contract}
Six terminal candidates are frozen before outcomes. Both memories use the same writer model; eligibility requires valid ReasoningBank items, paired embedding cosine at least 0.75, an independent verifier judging actionable guidance equivalent with no unique task-solving fact, and no literal benchmark-reference leakage. Four candidates pass: tasks $\{387,167,23,388\}$ with cosines 0.931, 0.902, 0.906, and 0.904. The same four remain eligible if the cosine floor is raised from 0.75 to 0.85 or 0.90; this is a zero-call re-thresholding of archived pre-outcome values, not an alternative-verifier run. Task 385 is excluded for repeating ``Cally''; task 164 repeats ``Dry'' and ``Uneven color.'' All exclusions are pre-outcome.

\subsection{Post-hoc operational-equivalence stress test}
The original information-equivalence gate remains the preregistered R4 admission rule; the following diagnostics do not retroactively change eligibility. An archived stronger-verifier audit judges none of the four admitted R4 pairs fully equivalent under one semantic reviewer and only one of four fully equivalent under a second, with the residual disagreements involving concrete strategy emphasis. A deterministic lexical audit points in the same direction without a new model call: failure-derived memory contains more search/filter language and more strong-directive language in all 6/6 R4 writer pairs; among the 24 R6 writer pairs it contains more failure/caution language in 23/24, more verification language in 20/24, and more strong-directive language in 20/24. These diagnostics make the limitation operational: the success/failure ReasoningBank intervention changes writer mode together with pragmatic register and strategy emphasis. They therefore weaken a provenance-only interpretation of L1 and increase the value of the byte-identical L2 intervention; they do not change either frozen L1 outcome.

\subsection{Terminal runtime repair boundary}
A condition-blind projection removes repeated global menu blocks, blank lines, and exact duplicate lines without consulting references or outcomes. An initial context-allocation support failure is discarded before scientific admission; the clean restart changes only context capacity. Its largest prompt uses 5,209 of 12,288 input tokens; executor, temperature, seed schedule, pair set, scorer, effect floor, permutation seed, and significance rule remain frozen.

\subsection{Terminal statistical rule}
Each qualified pair receives six paired seeds. The statistic is task-level terminal success, success-derived minus failure-derived. Support requires $\Delta\geq0.15$ and one-sided $p<0.05$; counterevidence requires $\Delta\leq-0.15$ and its one-sided $p<0.05$, using 100,000 frozen-seed permutations. Observed $\Delta=+0.1667$ with $p_S=0.0785$. Four independent task units imply only 16 exact sign assignments and a minimum one-sided value $1/16=0.0625$, so this endpoint cannot satisfy the frozen $0.05$ rule.

\subsection{Early-action cohort and L1 writer-mode contrast}
The disjoint WebArena cohort excludes prior D2 task IDs. A candidate must come from a released successful trajectory and expose a visible single-click reference action at step two or later; eligible tasks are sorted by ID and the first 24 are frozen before counterfactual writing. A compact process trace preserves instruction and action sequence but removes terminal success. The same trace is passed through first-party ReasoningBank success- and failure-reflection prompts using Qwen2.5-32B at temperature zero, directly varying outcome-conditioned writer mode while holding the trace fixed.

Eligibility again requires valid memory items, cosine at least 0.75, verifier-equivalent actionable guidance, no unique task-solving fact, and no reference-action leakage. Twenty-three pairs pass; candidate 159 is excluded for an unparsable verifier output. All 23 enter downstream execution.

\subsection{Early-action statistical rule}
Each task receives six paired seeds per writer condition with identical future query and frozen browser state. The output is one visible click index; the released successful trajectory supplies the reference. Inference is task-level: the permutation test sign-flips each task's success-derived minus failure-derived reference-match difference over 100,000 draws, with the same $\pm0.15$ magnitude and one-sided $0.05$ gates. The 23 tasks yield 276 complete requests and zero support failures; mean difference is $-0.0942$, with $p_S=0.9841$ and $p_F=0.0792$. Paired actions differ in 21.0\% of seeds; hypothesis-direction error-consistent divergence is 0.7\%, versus 10.1\% in reverse.

\subsection{Exact-information support-capacity stop and reopen condition}
R5 holds the actionable memory core byte-identical and changes only structured provenance metadata. Its frozen contract requires ten unique review tasks, but the released shopping split contains only five; it therefore stops before any model call, with no scientific verdict and no post-audit target reduction. The strongest reopen path is source-faithful financial AgentDojo replication of an L1 writer-mode or L2 exact-information contrast when query-level artifacts or auditable exact-model access become available. Such replication addresses L3 transport, not identification by itself; any new bridge must be frozen before outcomes and cannot reuse R4/R6 for threshold rescue.

\subsection{Prospective L1 power-planning contract}
\label{app:power}
The current $|\Delta|\geq0.15$ rule is a practical-relevance floor, not evidence of adequate power. A future confirmatory L1 run must choose the number of independent task pairs before outcomes by simulation or exact enumeration of the planned task-level sign-flip test. The planning distribution must be frozen from an external pilot or other prespecified source; the target effect is $|\Delta|=0.15$; the directional randomization threshold remains 0.05; and the target prospective power is at least 80\%. Pair construction, task exclusions, seeds per condition, maximum task count, and the stopping rule must be frozen with that calculation. Repeated requests within a task may reduce task-level measurement noise but do not count as additional independent units.

For scale, a normal-approximation planning sensitivity for a paired task-level mean at $|\Delta|=0.15$ gives the following independent-task counts. These are design aids, not retroactive evidence, because the unknown task-level standard deviation must be fixed prospectively from an external pilot.
\begin{center}
\begin{tabular}{lrr}
\toprule
Planning SD & one-sided $\alpha=.05$, 80\% power & two-sided $\alpha=.05$, 80\% power \\
\midrule
0.20 & 11 & 14 \\
0.30 & 25 & 32 \\
0.40 & 44 & 56 \\
\bottomrule
\end{tabular}
\end{center}
Thus ``four tasks'' is not defended as an adequately powered design; even the optimistic SD=0.20 planning case points to a double-digit confirmatory cohort. R4 remains structurally non-decisive, and R6 is judged by its frozen effect and significance gates rather than post-hoc sample extension. We retain the preregistered directional randomization tests as the decision tests. A future confirmatory packet should additionally report a two-sided randomization sensitivity and a paired confidence interval; BCa intervals are descriptive rather than a substitute for the frozen randomization gate, and are not retrofitted onto the four-task R4 endpoint.

\subsection{Source-faithful L3 decision contract}
\label{app:l3}
An L3 transport test must recover the motivating financial-agent writer/checkpoint and original memory interface, freeze the same acquisition trajectory content and downstream schedule across conditions, and manipulate only the declared provenance channel before any outcome is observed. For an L1 source-native test, success- versus failure-conditioned writer mode may change the generated memory, after which the pre-outcome information-equivalence audit is applied and the paired memories are executed in the original runtime. For an L2 source-native test, actionable memory, retrieval membership, and retrieval order must remain fixed while only executor-visible provenance information changes (e.g., truthful provenance hidden versus exposed).

The current obstacles separate cleanly. \emph{Infrastructure debt} is the missing query-level ReasoningBank artifacts, auditable access to the reported Qwen 3.7 Flash writer and qwen3.7-text-embedding configuration, source checkpoint/prompt bindings, and provenance-bearing runtime fields needed to replay the original financial tasks. \emph{Conceptual requirements} remain even after those assets exist: the intervention channel must be declared before outcomes; L1 must retain a pre-outcome information-equivalence gate; L2 must keep actionable memory byte-identical; and any WebArena-to-financial transport claim must state its target-population assumptions rather than treating source replication as automatic generalization. A writer-output change plus downstream change supports a source-native writer-mode path; a byte-identical-memory metadata contrast with downstream change supports an L2 metadata path on that substrate; and adequate-support null evidence on both channels is substrate-bounded counterevidence rather than a general no-provenance theorem. No such source-faithful execution is claimed in the present paper.

\subsection{Fail-closed R19 stop and non-reuse rule}
A later L2 compatibility experiment prospectively freezes 35 independent tasks and 140 terminal episodes under a two-sided task-level sign-flip primary analysis. It reaches 29/140 durable episodes: seven tasks are complete and one episode of an eighth task is complete. The next reset then times out before STARTED; the single exact pre-exposure retry allowed by the frozen contract also times out before STARTED. Because the retry budget is exhausted, the confirmatory attempt stops. No task-level deltas, effect mean, $p$-value, interval, or subgroup result is computed from the partial prefix for manuscript claims. The prefix may not be pooled into a later cohort and may not select a new threshold, endpoint, model, controller, or statistical rule.

After excluding every template exposed by the stopped attempts, the same frozen asset contains 27 fully unexposed templates. This is inventory, not an automatically authorized replacement cohort. For reference, the preregistered planning sensitivity at task-level SD $0.30$ requires about 32 independent tasks for a two-sided 80\% target at $|\Delta|=0.15$; $n=27$ is below that reference. Sample scarcity is therefore not a reason to lower the gate.

\subsection{Fresh R53--R57 L2 execution contract and receipts}
The completed L2 experiment is a new lineage rather than a continuation of R19. R53 executes all 350 tasks in the frozen OSInteraction training split under a single exactly-once source-memory attempt, producing 176 success-derived and 174 failure-derived memories. Its source receipt records 350/350 completed units and zero external-provider calls. R54 evaluates retrieval support before any validation environment reset. Historical validation signatures are excluded; among 108 fresh exact skill-signature clusters, 106 have eligible native retrieval. Hash ranking freezes 32 primary and eight utilization representatives before the first validation treatment.

R55 then executes 40 utilization arms (eight clusters times five memory surfaces) with started/completed durability ledgers. It completes 40/40 without failure. The frozen promotion rule uses the first executable action only: true memory is specifically divergent from both no-memory and null-memory in 5/8 clusters, versus 3/8 null-memory versus no-memory divergences, so the first stage passes. Terminal success in R55 is diagnostic and is not used to select the primary sample.

R56 freezes the A/B schedule before its first treatment. The R39 exact-information adapter renders each selected retrieval into A (explicit source-outcome field masked) and B (the same rows/content/order plus truthful boolean source outcome). The adapter audit requires actionable-content identity and verifies that \texttt{source\_outcome\_success} is the only executor-visible field added in B; similarity, Q values, scores, IDs, and source task IDs are hidden. Memory text itself is not scrubbed of implicit outcome clues, so A is a masked-field baseline rather than a no-provenance-information condition. All 64 arms complete with zero execution failures. Complete-only analysis yields A=15/32, B=16/32, observed $\Delta=0.03125$, one discordant pair, exact two-sided sign-test $p=1.0$, and the frozen 100,000-resample paired bootstrap interval $[0,0.09375]$. The observed effect does not reach the 0.15 threshold; no equivalence claim is implied.

\subsection{R59--R62 executor-backbone replication}
The second-backbone experiment is prospectively frozen after R57 and before any Llama validation outcome. The local Meta-Llama-3.1-8B-Instruct artifact is content-addressed, and a 3/3 source-side native-parser qualification passes without opening validation or terminal outcomes. The replication then reuses the exact R54 source bank and frozen retrieval content/order, the same eight utilization and 32 primary cluster IDs, the same R39 exact-information A/B renderer, and the same evaluator/randomization/analysis; only the executor backbone changes.

R60 executes 40 Llama utilization arms and completes 40/40 without failure. The memory-specific first-action criterion holds in 8/8 clusters, while null-memory versus no-memory differs in 5/8, so the frozen promotion gate passes before any primary outcome. R61 then executes the 64 frozen A/B arms exactly once. Complete-only analysis gives an observed $\Delta=0$, two B-only and two A-only successes, four discordant pairs, exact two-sided sign-test $p=1.0$, and the frozen 100,000-resample paired bootstrap interval $[-0.125,0.125]$. The observed effect does not reach the 0.15 threshold. R62 adjudicates Qwen and Llama separately, without statistical pooling. This is executor-only robustness within one shared source-bank/retrieval/task/renderer substrate, not an independent pipeline replication or universal zero-effect claim.

\subsection{R66 post-confirmatory sparse-discordance and isolation audit}
R66 changes no experimental outcome, threshold, frozen analysis, or execution authority. First, it adds a conservative paired-risk-difference audit for the sparse binary discordances. Let $p_{10}=P(B=1,A=0)$ and $p_{01}=P(B=0,A=1)$, so the paired risk difference is $p_{10}-p_{01}$. We form separate exact two-sided 97.5\% Clopper--Pearson intervals for $p_{10}$ and $p_{01}$ and combine them by Bonferroni as $[L_{10}-U_{01},U_{10}-L_{01}]$, guaranteeing at least 95\% simultaneous coverage. This yields Qwen $[-0.128,0.183]$ and Llama $[-0.225,0.225]$. Consequently neither model establishes $\pm0.15$ equivalence and neither conservatively excludes a $+0.15$ population effect. These intervals are a post-confirmatory wording audit and do not replace the preregistered bootstrap/sign test.

Second, R66 audits arm isolation at the bound implementation. The shared \texttt{run\_arm\_exact} creates a fresh task and session per arm, invokes \texttt{task.reset(session)} before the inference loop, and invokes \texttt{task.release()} in a \texttt{finally} block. At pinned MemRL revision \texttt{c1b322ca}, OSInteraction reset constructs a new \texttt{OSInteractionContainer}; its container constructor launches a fresh Docker container with \texttt{remove=True}, while normal completion and release terminate the container. Thus one arm is not intended to inherit task-visible users, files, or processes from the preceding arm. The experiment holds the reset initial environment-generation mechanism fixed; once treatment changes an action, later observations are permitted to differ as treatment-induced mediators. The original ledgers do not contain a dedicated per-arm \texttt{reset\_success} timestamp bit, so the audit is code-level reset/container evidence plus 64/64 completed-arm evidence rather than a newly manufactured runtime receipt.

\subsection{Future transport and governance gate}
\label{app:fresh}
L3 transport or PSMG-efficacy evaluation remains a new scientific object and opens only after its treatment-independent components are frozen. For PSMG efficacy, the provenance-sensitive controller, strongest same-information provenance-free comparator, calibration data, decision thresholds, and action mapping must be executable and content-addressed before target outcomes. For L3, the motivating source runtime and artifacts must be recovered. The completed Qwen and Llama L2 results remain fixed regardless of future experiment direction and may not be reinterpreted as failed attempts to be rescued.
